\documentclass[UTF8,a4paper,11pt]{ctexart}

\usepackage[margin=2.4cm]{geometry}
\usepackage{amsmath}
\usepackage{amssymb}
\usepackage{booktabs}
\usepackage{enumitem}
\usepackage{hyperref}
\usepackage{xcolor}
\usepackage{listings}
\usepackage{fancyhdr}

\setlength{\headheight}{13.6pt}
\pagestyle{fancy}
\fancyhf{}
\fancyhead[L]{\small\kaishu\nouppercase{\leftmark}}
\fancyfoot[C]{\thepage}
\renewcommand{\headrulewidth}{0pt}
\renewcommand{\footrulewidth}{0pt}
\fancypagestyle{plain}{
  \fancyhf{}
  \fancyfoot[C]{\thepage}
  \renewcommand{\headrulewidth}{0pt}
  \renewcommand{\footrulewidth}{0pt}
}

\hypersetup{
  colorlinks=true,
  linkcolor=blue!60!black,
  urlcolor=blue!60!black,
  citecolor=blue!60!black
}

\lstdefinestyle{code}{
  basicstyle=\ttfamily\small,
  breaklines=true,
  frame=single,
  columns=fullflexible,
  keepspaces=true
}
\lstset{style=code}

\newcommand{\stage}[1]{\text{\texttt{#1}}}
\newcommand{\nextstage}{\;\longrightarrow\;}

\title{基于 SysY 的多模式编译器设计与实现报告}
\author{（朱宏泽 / 2400012623）\\
\small 项目仓库：\href{https://github.com/glazecats/koopa-compiler}{glazecats/koopa-compiler}}
\date{\today}

\begin{document}

\maketitle

\begin{abstract}
本项目实现了一个以 SysY/类 C 语言为基础的课程编译器。与只追求单一路径
``能编译'' 不同，本项目将编译流程拆分为三个面向不同目标的模式：
\texttt{-riscv} 面向基础正确编译与 RISC-V 汇编输出，
\texttt{-perf} 面向性能优化与分配实验，
\texttt{-extension} 面向语言特性扩展。三条模式共享前端、语义分析、分层
IR、SSA 与后端基础设施，但各自采用不同的保守边界和测试策略。项目的核心
特色包括 canonical IR 与 lower IR 的分层设计、ValueSSA/MemorySSA 风格的中端
结构、family-aware 的 SSA 分配实验，以及以函数值、闭包和 callable-object IR
为代表的语言扩展主线。开发过程中采用快速迭代的 vibe coding 工作流，但每个
迭代都由 verifier、dump regression、runtime regression 和全量测试闭环约束；
也就是说，快速实现必须始终服从可验证的工程纪律。
\end{abstract}

\section{项目目标与总体思路}

本项目的目标是实现一条完整、可验证、可扩展的编译器主线。输入语言以 SysY
风格的类 C 语言为基础，输出方向以 RISC-V 汇编为主，同时保留面向中间表示
观察、优化实验和语言扩展的多条工程路径。

从整体结构看，编译流程可以概括为：

\[
\begin{aligned}
\stage{source}
&\nextstage \stage{lexer}
\nextstage \stage{parser}
\nextstage \stage{AST}
\nextstage \stage{semantic} \\
&\nextstage \stage{canonical IR}
\nextstage \stage{lower IR}
\nextstage \stage{ValueSSA / MemorySSA} \\
&\nextstage \stage{machine backend}
\nextstage \stage{RISC-V asm}
\end{aligned}
\]

这个项目最重要的设计选择，是没有把所有能力塞进一条难以控制的“大管线”，
而是按目标拆成三个模式：

\begin{enumerate}[label=\arabic*.]
  \item \texttt{-riscv}：面向课程基础目标，强调正确翻译、稳定输出和后端连通。
  \item \texttt{-perf}：面向性能目标，在正确性闭环下进行优化和寄存器分配实验。
  \item \texttt{-extension}：面向语言扩展，用保守 pipeline 接纳更复杂的语义特性。
\end{enumerate}

三种模式不是三套互不相干的编译器，而是同一个工程在不同约束下的三种使用方式。
基础输出、性能优化和语言扩展放在一起，使项目更接近真实编译器工程，
而不只是完成单一评测目标的最小实现。

\subsection{评价目标与边界}

除了三个主模式外，driver 还提供了一个可叠加的严格检查开关：
\begin{lstlisting}[language={}]
--enforce-all-paths-return-check
\end{lstlisting}
它不是第四种编译模式，而是放在 \texttt{-riscv}、\texttt{-perf}、
\texttt{-extension} 前面的语义审计选项，
用于更严格地检查非 \texttt{void} 函数是否在所有控制流路径上返回。它也会提前拦截
一类静态可判定的非终止/不可达返回问题：例如明显会卡在 \texttt{while(1)} 或
\texttt{for(;;)} 中的路径，不能靠循环后面不可达的 \texttt{return} 来伪装成
“所有路径都返回”。默认编译路径为了兼容课程提交和较宽松的 SysY 测试习惯，会有意
绕开这一严格 gate；当需要审计控制流返回性质时，再显式打开该选项。

因此，本项目的效果评价也按目标拆开：

\begin{itemize}
  \item \texttt{-riscv}：重点看基础程序能否稳定通过语义检查、lowering 和后端，最终生成可汇编的 RISC-V 文本；
  \item \texttt{-perf}：重点看优化是否带来真实运行收益，并通过 A/B、compile-time 和正确性回归共同决定是否保留；
  \item \texttt{-extension}：重点看复杂语言特性能否被诚实地分阶段支持，支持的形状要有 runtime 或 dump regression，暂不支持的形状要明确 gate；
  \item strict return-check：重点看控制流返回分析是否足够保守、可解释，并避免把合法程序误判为缺失返回。
\end{itemize}

这也说明了项目的一个 non-goal：它并不试图一次性实现完整工业级 C 编译器或完整一等函数
运行时，而是在课程编译目标之上，逐步把 IR、SSA、后端和扩展语义做成可验证、可继续生长
的工程。换句话说，暂未开放的能力应当由明确的 gate 拒绝，而不是在边界不清的位置
生成不可靠的中间表示。

\subsection{工程目录结构}

从仓库结构上看，项目并不是把所有实现压在一个目录里，而是按编译阶段和工程职责拆分：

\begin{itemize}
  \item \texttt{src/} 与 \texttt{include/}：编译器实现与公共头文件，包括前端、IR、SSA、后端和 runtime/observe 相关模块；
  \item \texttt{tests/}：按 lexer、parser、semantic、IR、lower IR、SSA、machine、driver、extension runtime 等层次组织的回归测试；
  \item \texttt{docs/}：工程记忆、下一步计划、阶段设计文档和优化/语言特性路线；
  \item \texttt{lesson/}：由 review/lesson agent 精炼出的项目说明材料，用于给报告和后续维护提供压缩上下文；
  \item \texttt{tools/}：dump、profile、diagnostic 等辅助工具，用于观察中间阶段；
  \item \texttt{reports/}：课程报告源文件与生成的 PDF。
\end{itemize}

这样的目录结构也对应了项目的协作方式：实现、测试、设计记录、精炼 lesson 和课程报告
各自有位置，不需要把所有上下文集中到单一总文档中。

\section{公共基础设施：分层 IR 与验证体系}

\subsection{前端与语义分析}

前端部分包括词法分析、语法分析、AST 构造和语义分析。词法分析负责将字符流
切分成 token；语法分析将 token 组织成声明、语句和表达式；AST 作为 source-near
的结构化表示，保留后续语义检查和 lowering 所需的信息。语义分析则负责名字
解析、作用域、类型检查、函数签名、控制流返回约束等规则。

项目后期加入了不少扩展特性，例如 \texttt{defer}、函数值、闭包、结构体、
保守浮点支持等，因此语义分析不只负责基础名字解析。它还承担了模式边界
控制：哪些语法只在 \texttt{-extension} 下开放，哪些组合仍然应该被拒绝，必须
在语义层给出明确诊断。

虽然这些前端步骤在课程编译器中相对标准，但它们仍然是后续 IR 正确性的入口。
例如下面这段源码：

\begin{lstlisting}[language=C]
int add(int x, int y) {
  return x + y;
}
\end{lstlisting}

词法分析会先把字符流切成关键字、标识符、标点和运算符：

\begin{lstlisting}[language={}]
INT IDENT(add) LPAREN INT IDENT(x) COMMA INT IDENT(y) RPAREN
LBRACE RETURN IDENT(x) PLUS IDENT(y) SEMICOLON RBRACE
\end{lstlisting}

语法分析和 AST 构造再把这些 token 组织成“函数定义、参数列表、返回语句、二元表达式”
这样的树形结构：

\begin{lstlisting}[language={}]
FunctionDef name=add return=int
  Param x:int
  Param y:int
  Body
    Return
      Binary(+)
        Identifier(x)
        Identifier(y)
\end{lstlisting}

语义分析则在 AST 上补充更“看不见”的约束：\texttt{x} 和 \texttt{y} 是否在当前
作用域可见，\texttt{x + y} 的类型是否为 \texttt{int}，返回表达式是否匹配函数返回
类型，以及非 \texttt{void} 函数是否满足当前模式要求的返回路径规则。换句话说，
lexer/parser/AST 给出程序的结构，semantic 判断该结构是否满足语言规则。

语义层还负责把一些错误尽早挡在 IR 之前。例如：

\begin{lstlisting}[language=C]
int bad(int x) {
  if (x > 0) return x;
}
\end{lstlisting}

在默认课程兼容路径下，这类程序可以按较宽松规则处理；但打开严格返回检查开关后，
semantic 会指出非 \texttt{void} 函数存在未返回路径。类似地，若程序在非 \texttt{-extension}
模式下使用闭包或函数值，或者在 \texttt{-extension} 下写出尚未开放的函数值组合，
也会在语义阶段被明确拒绝，而不是生成语义契约已经不可靠的中间表示。这个前端的目标
不是增加表面复杂度，而是把模式边界和错误诊断尽量前置：通过的程序可以交给后续
verifier 和 lowering，未通过的程序也要尽量说明拒绝原因。

\subsection{canonical IR}

canonical IR 是语义分析之后的第一层稳定中间表示。它是 block-based、非 SSA
的 IR，保留较多接近源码的控制流结构，适合做早期规范化和轻量清理。

这一层的主要职责是：

\begin{itemize}
  \item 将 AST 中的表达式、语句、控制流转成统一 IR；
  \item 用基本块和 terminator 显式表达控制流；
  \item 通过 verifier 检查 block、temp、local、call、return 等结构契约；
  \item 运行常量折叠、copy propagation、CFG cleanup 等轻量 pass。
\end{itemize}

\subsection{lower IR}

lower IR 的核心思想是显式区分 value 和 storage。canonical IR 中的 local/global
在 lower IR 中不再随意作为普通 value 出现，而是通过 \texttt{load\_*} 和
\texttt{store\_*} 指令访问。

这一层的意义在于：它给后续 SSA construction 提供了更干净的输入。ValueSSA
只需要处理 value flow，而不必同时承担完整 memory SSA 的语义责任。
这是一种很“工程化”的分层：先把边界说清楚，再往后做复杂分析。

\subsection{ValueSSA 与 MemorySSA}

ValueSSA 将 lower IR 的值流转换为 SSA 形式，引入显式的 SSA value、def-use
关系和 phi 节点。这里的 SSA value 就是静态单赋值意义上的值名：每个
\texttt{ssa.*} 只定义一次，后续使用都通过 def-use 关系追踪。它主要处理值流，
而不是完整内存状态。MemorySSA 则作为后续补充，用于表达和优化内存版本、
load/store 关系和内存依赖。

这种拆分使得项目可以分别推进：

\begin{itemize}
  \item ValueSSA 上的常量折叠、冗余二元表达式消除、CFG 简化、dead value DCE；
  \item MemorySSA 上的 load forwarding、store cleanup、memory CSE；
  \item allocator 和 machine backend 所需的 liveness、interference、allocation prep。
\end{itemize}

\subsection{分层 IR 示例}

下面用一个很小的条件表达式说明三层 IR 的差异。源码本身并不复杂，但它已经
包含控制流汇合：

\begin{lstlisting}[language=C]
int choose(int a, int b, int c) {
  return c ? a + 1 : b + 2;
}
\end{lstlisting}

在 canonical IR 中，条件表达式被拆成显式基本块。两个分支分别计算结果，并汇合到
同一个临时值 \texttt{tmp.0}：

\begin{lstlisting}[language={}]
func choose(a.0, b.1, c.2) {
  bb.0:
    br c.2, bb.1, bb.2
  bb.1:
    tmp.1 = add a.0, 1
    tmp.0 = mov tmp.1
    jmp bb.3
  bb.2:
    tmp.2 = add b.1, 2
    tmp.0 = mov tmp.2
    jmp bb.3
  bb.3:
    ret tmp.0
}
\end{lstlisting}

lower IR 则把 storage 边界显式化。参数和局部变量不再被当作普通值随意消费，
而是先通过 \texttt{load\_local} 读出；这里没有局部变量写回，所以不会出现
\texttt{store\_local}：

\begin{lstlisting}[language={}]
func choose(a.0, b.1, c.2) {
  bb.0:
    tmp.3 = load_local c.2
    br tmp.3, bb.1, bb.2
  bb.1:
    tmp.4 = load_local a.0
    tmp.1 = add tmp.4, 1
    tmp.0 = mov tmp.1
    jmp bb.3
  bb.2:
    tmp.5 = load_local b.1
    tmp.2 = add tmp.5, 2
    tmp.0 = mov tmp.2
    jmp bb.3
  bb.3:
    ret tmp.0
}
\end{lstlisting}

ValueSSA 再把 value flow 改写成静态单赋值形式。两个分支中产生的值分别是
\texttt{ssa.3} 和 \texttt{ssa.6}，汇合块用 \texttt{phi} 根据控制流前驱选择
正确来源：

\begin{lstlisting}[language={}]
func choose(a.0, b.1, c.2) {
  bb.0:
    ssa.0 = load_local c.2
    br ssa.0, bb.1, bb.2
  bb.1:
    ssa.1 = load_local a.0
    ssa.2 = add ssa.1, 1
    ssa.3 = mov ssa.2
    jmp bb.3
  bb.2:
    ssa.4 = load_local b.1
    ssa.5 = add ssa.4, 2
    ssa.6 = mov ssa.5
    jmp bb.3
  bb.3:
    ssa.7 = phi [bb.1: ssa.3], [bb.2: ssa.6]
    ret ssa.7
}
\end{lstlisting}

这个例子很小，但能看出分层的核心：canonical IR 方便从 AST 生成和做早期清理；
lower IR 把 value/storage 说清楚；ValueSSA 再用静态单赋值和 \texttt{phi} 把
跨控制流的值选择显式化。每层只解决本层的结构问题，可以减少后续阶段的隐式假设和调试成本。

\section{\texttt{-riscv} 模式：基础编译与 RISC-V 输出}

\texttt{-riscv} 模式是项目的基础交付模式。它面向课程编译器的核心目标：将
SysY/类 C 程序正确编译为 RISC-V 汇编，并尽量保持输出可被标准工具链接受。

\subsection{设计目标}

\texttt{-riscv} 模式强调三个关键词：

\begin{enumerate}[label=\arabic*.]
  \item 正确性：源程序语义必须被保守、稳定地翻译。
  \item 可汇编性：输出文本应能通过 RISC-V 汇编器/链接器检查。
  \item 可维护性：后端每一层都有自己的 artifact、summary 和测试入口。
\end{enumerate}

在该模式下，语言能力主要围绕课程基础需求展开，不默认接受所有
\texttt{-extension} 特性。这避免了扩展语义对课程基础路径造成额外风险。

\subsection{后端分层}

后端主线大致可以拆成：

\[
\begin{aligned}
\stage{machine\_ir}
&\nextstage \stage{machine\_select}
\nextstage \stage{machine\_layout} \\
&\nextstage \stage{machine\_emit}
\nextstage \stage{machine\_encode}
\nextstage \stage{RISC-V text}
\end{aligned}
\]

其中，machine IR 承接 ValueSSA/allocator 之后的 machine-facing 表示；
machine select 将较抽象的操作选择成更接近目标机器的 selected op；
layout/emit/encode 则逐步处理块布局、标签、目标 profile 和最终 RISC-V 文本输出。

除了这条主交付路径外，项目还建立了 object/runtime/observe 相关的支撑层，例如
bytes、object、reloc、ELF、image、runtime、trace、journal 等。它们在本文中不作为
独立主线展开，但很能说明后端的工程取向：后端不是一个把字符串拼成汇编的
``printf 汇编'' 文件，而是一组带 artifact、summary、verifier 和测试入口的阶段。
这样做的好处是，某一层出问题时可以从中间产物定位；某一层需要替换时，也不必把整个
后端拆成一地零件。

\subsection{后端落地示例}

沿用上一节的 \texttt{choose} 例子：

\begin{lstlisting}[language=C]
int choose(int a, int b, int c) {
  return c ? a + 1 : b + 2;
}
\end{lstlisting}

ValueSSA 中的 \texttt{phi} 表达了“从哪个分支来的值”。进入后端后，这个例子会被进一步
整理成两个直接返回的分支，因此 machine-facing 形状已经不再保留显式汇合块：

\begin{lstlisting}[language={}]
machine_select
function choose params=3 locals=3 spills=0
  bb.0:
    reg.0 = load_local local.2
    br reg.0, bb.1, bb.2
  bb.1:
    reg.0 = load_local local.0
    reg.0 = alui.0 reg.0, 1
    ret reg.0
  bb.2:
    reg.0 = load_local local.1
    reg.0 = alui.0 reg.0, 2
    ret reg.0
\end{lstlisting}

最终输出的 RISC-V 文本则体现为参数落栈、条件分支、两个返回路径和栈帧恢复：

\begin{lstlisting}[language={}]
choose:
  addi sp, sp, -16
  sw a0, 0(sp)
  sw a1, 4(sp)
  sw a2, 8(sp)
.Lchoose_0:
  lw a0, 8(sp)
  bne a0, zero, .Lchoose_2
.Lchoose_1:
  lw a0, 4(sp)
  addi a0, a0, 2
  addi sp, sp, 16
  ret
.Lchoose_2:
  lw a0, 0(sp)
  addi a0, a0, 1
  addi sp, sp, 16
  ret
\end{lstlisting}

这个片段展示了 \texttt{-riscv} 模式的核心效果：前面 IR 中较抽象的控制流和值选择，
在后端被转成目标机器可接受的标签、分支和寄存器/栈操作。中间形状可以变化，但每一步
都有 dump 可以观察，因此比只检查最终汇编文本更容易定位问题。

\subsection{效果}

\texttt{-riscv} 模式最终希望达到的效果是：

\begin{itemize}
  \item 能完成基础 SysY 程序的前端检查、中端 lowering 和后端输出；
  \item 生成 assembler-accepted 的 RISC-V preview asm；
  \item 对函数调用、全局变量、数组/间接内存、控制流、返回值等基础语义保持正确；
  \item 在后端 cleanup 中避免危险 peephole，例如跨 call 错误复用 caller-clobbered register。
\end{itemize}

\section{\texttt{-perf} 模式：优化、SSA pass 与分配实验}

\texttt{-perf} 模式面向性能优化。与 \texttt{-riscv} 相比，它更关注运行时间、
热点程序和优化收益，但仍然以正确性为前提。

\subsection{优化工作流}

项目的性能优化工作流遵循两个紧凑原则：优化候选先来自热点或结构证据，是否保留再由
A/B、正确性和编译时间共同决定。也就是说，优化不是看到汇编里某个形状不理想就直接改，
而是结合：

\begin{itemize}
  \item 源程序的循环、递归和调用结构；
  \item 生成汇编中的重复计算、地址链、调用开销；
  \item 实际运行时间和热点 witness；
  \item 编译时间变化和正确性回归风险。
\end{itemize}

这种流程在课程项目里略显“重”，但很有价值。因为优化很容易产生误导性信号：
静态代码形状变短，并不必然意味着运行时间变短，更不必然意味着语义仍然正确。
因此优化候选必须经过运行结果、性能数据和回归测试共同确认。

\subsection{多层优化主线}

\texttt{-perf} 的主要优化发生在 ValueSSA 及其相邻层，但并不只发生在 SSA。项目更接近
一组分层的优化实验场：前面的 IR 层适合做语义更清楚的规范化，中间 SSA 层适合做
def-use 和控制流优化，后端层则适合做目标相关的局部清理。粗略分布如下：

\begin{itemize}
  \item canonical IR：早期 constant folding、copy propagation、CFG cleanup、死临时值清理；
  \item lower IR：value/storage boundary normalization，为后续 SSA 和 memory analysis 减少歧义；
  \item ValueSSA：trivial value simplification、algebraic identity simplification、constant folding、redundant binary elimination、CFG simplification、dead value DCE、tiny helper inlining、LICM/induction/address-carrier 相关 hotspot pass；
  \item MemorySSA：local/global load forwarding、redundant/dead store cleanup、local/global slot scalar-replacement 风格实验；
  \item machine IR / machine select：slot-aware cleanup、跨块 alias/copy propagation、基于 live-out 的 dead-def elimination、call barrier 下的 register/spill-slot 区别处理；
  \item final RISC-V text：少量 target-specific peepholes，例如 tail-return、repeated-address pattern、stack store/reload 等，但这层必须最谨慎。
\end{itemize}

这些 pass 的共同特点是：尽量把优化放在语义信息还足够丰富的层上，优先做可验证的小步
变换；最终文本 peephole 只处理很局部、目标相关、并且有回归保护的形状。经验上，越靠近
最终汇编，可用的语义信息越少，因此文本级优化应当限制在非常具体、可回归验证的模式上。

\subsection{family-aware 的 SSA 分配器}

\texttt{-perf} 模式中另一个有特色的部分是 ValueSSA allocator。它的基本骨架仍然接近
传统图染色分配器：先根据 liveness 建 interference graph，再把低压力节点
simplify/remove 到一个栈式顺序里，最后按逆序 select color；如果某个 value 在逆序染色时
没有可用颜色，就进入 spill / rewrite 路线。

项目里的实现没有推翻这个骨架，而是在骨架上加了一层 family-aware policy。朴素实现通常
把每个 SSA value 作为相对独立的分配对象处理，逆序染色时能分配则分配，不能分配则 spill；
这里会额外问：
哪些 value 来自 \texttt{mov}/\texttt{phi}/copy affinity，应该尽量被视为同一组？
如果它们想共色，当前 blocker 是不是可以被安全挪开？如果某个值先被当成 spill，
后面有没有机会 recover？

更准确的主线可以概括为：

\begin{lstlisting}[language={}]
SSA values
  -> interference / liveness / affinity
  -> coalesce pair analysis
  -> coalesced family / root
  -> family-aware simplify / freeze / spill plan
  -> reverse select coloring
       with preferred-color / targeted-eviction / retry-recovery
  -> spill slot / SSA rewrite / canonicalize / layout report
\end{lstlisting}

这个流程包含多个阶段。第一步 liveness/interference 先回答“哪些 SSA value 的生命周期
重叠，因此不能放在同一个寄存器颜色里”；affinity/coalesce analysis 再从 \texttt{mov}、
\texttt{phi} 和 copy-like value flow 里找出“哪些值最好能共色”。随后 planner 并不是直接
染色，而是先根据 pressure、move 关系和 spill 风险生成 removal plan；真正的颜色选择发生在
reverse select coloring 阶段。最后，如果确实 spill，就进入 spill-slot placement、SSA
rewrite 和后续 canonicalize。换句话说，报告里的 allocator dump 是一组阶段化诊断信息，
包括关键决策点的状态快照和 allocation trace。

其中，\texttt{family} 表示当前被视为同一 coalesced class 的一组值，\texttt{root} 是
family 的代表。当前实现还不是教材里那种完全 merged-node 的 Briggs/Chaitin allocator：
最终结果仍然落回每个 \texttt{ssa.N} 的 color/spill placement；但 planner 和 selector
已经会利用 family 信息来调整过程。

这个 family 更接近一种 soft coalescing 视角。实现上先用 union-find 风格的 root 和
class-size 数组记录哪些 value 属于同一组，但不会把整组立刻压成一个真正的 super-node。
remove 时仍然是一项一项移出，例如先移出 \texttt{ssa.3}，
再移出 \texttt{ssa.1}；即使 root 对应的 value 先被移走，root 也仍然作为 family id 存在。
因此实现里会把三件事分开维护：family identity、当前还 active 的 family 成员数，以及
reverse select 时已经染色的 family 成员。plan 完成后还会把动态变化后的 family 写入：

\begin{lstlisting}[language={}]
final_runtime_coalesce_roots
final_runtime_coalesce_class_sizes
\end{lstlisting}

select 阶段再通过这些 root/class 信息寻找 family preferred color。

实现上比较关键的概念不是单纯拿来展示，而是会直接进入
allocator 的决策：

\begin{itemize}
  \item \texttt{degree} / \texttt{effective\_degree}：planner 用它判断寄存器压力。原始 \texttt{degree} 来自 interference graph；\texttt{effective\_degree} 会扣掉 family 内部已经可解释的压力，从而避免把“本来希望共色的一组值”误判成更难分配；
  \item \texttt{move\_related} / \texttt{effective\_move\_related}：planner 和 move-family worklist 用它决定是继续尝试 coalesce，还是 freeze 掉 move 关系再进入 simplify/spill。前者看原始 move 边，后者只保留 family 外部仍然需要处理的 move 压力；
  \item \texttt{work\_class} / \texttt{move\_work\_class} / \texttt{phase}：这些字段决定候选进入哪个队列或阶段，例如 simplify、freeze、coalesce-ready、spill。也就是说，它们先决定“现在应该按哪类策略处理”，后面的 \texttt{priority} / \texttt{move\_work\_priority} / \texttt{spill\_cost} 再在对应策略里排序；
  \item \texttt{preferred color}：select 阶段的颜色偏好。普通图染色只要找任意可用颜色即可；这里会优先尝试 partner/family 已经使用的颜色，以减少后续 move；
  \item \texttt{targeted eviction}：当 preferred color 只被少数 blocker 挡住时，selector 会尝试定向挪开 blocker，再保住当前 value 的 preferred color；如果不安全，才退回普通可用颜色或 spill/retry 路线；
  \item \texttt{spill\_intended}、\texttt{spill\_confirmed}、\texttt{spill\_recovered}：这些字段把 planner 的预测和 selector 的最终结果分开。它们让 rewrite 阶段只处理真正 confirmed spill，也让 regression 能检查某个高风险值是否被 optimistic coloring 或 retry 成功救回。
\end{itemize}

\texttt{priority} 是这里最容易误读的字段。它不是一个全局绝对分数，也不是“越大就一定越
应该留在寄存器里”。当前实现大致把 affinity、interference degree、live block count 和
use count 等信号合成一个 base heuristic；真正比较时还要看 phase：

\begin{itemize}
  \item simplify 阶段优先选择已经低压、可安全摘除的 value。这里较高的 \texttt{priority} 或 move-work priority 更多表示“这个安全候选对 family/move 调度更有影响”，而不是“它更危险”；
  \item coalesce-ready 阶段会把 partner weight、mutual-best、family size、external affinity 等信号加进 \texttt{move\_work\_priority}，用来判断是否值得先合并 family，避免太早 remove 导致错过共色机会；
  \item freeze 阶段处理的是“move 关系还在，但暂时不继续 coalesce”的 family。freeze 会放弃这部分 move 优化偏好，但 interference 约束仍然保留；之后该 family/value 更容易回到普通 simplify 路线；
  \item spill 阶段不只看 \texttt{priority}，还会优先比较 \texttt{spill\_cost}、rematerializable、split-child、family pressure 等信号。\texttt{spill-remove} 也只是 planner 的高风险预测，不等于最终 confirmed spill。
\end{itemize}

因此，\texttt{priority=40} 的 \texttt{simplify-remove} value 和 \texttt{priority=30} 的
\texttt{spill-remove} value 不能直接跨类比较。前者是在“安全摘除”队列里更值得安排，
后者是在“可能 spill”队列里风险较高；二者表达的是不同阶段的调度含义。

下面的例子不是从某段 C 源码自然编译出来的，而是 allocator 层的 regression witness。
测试直接构造一个很小的 interference / coalescing 配置，用来锁定“高风险值在 select 阶段
如何通过 eviction/recolor 保住颜色”的策略。这样做的好处是变量少、压力关系清楚；
缺点是它不像源码例子那样直观，所以需要先说明设定：

\begin{lstlisting}[language={}]
color budget:
  3 colors

coalesced families:
  family A: {ssa.1, ssa.3}, root = ssa.1
  family B: {ssa.2, ssa.5}, root = ssa.2

selected interference:
  ssa.0 conflicts with ssa.1, ssa.2, ssa.6
  ssa.1 conflicts with ssa.4
  ssa.3 conflicts with ssa.4, ssa.5
  ssa.4 conflicts with ssa.5, ssa.6
  ssa.5 conflicts with ssa.6

removal plan:
  ssa.0: spill-remove, priority = 30
  ssa.3: simplify-remove, priority = 8
  ssa.1: simplify-remove, priority = 10
  ssa.6: simplify-remove, priority = 40
  ssa.5: simplify-remove, priority = 6
  ssa.2: simplify-remove, priority = 10
  ssa.4: simplify-remove, priority = 5
\end{lstlisting}

也就是说，这个 witness 故意制造了一个局面：\texttt{ssa.0} 是一个
\texttt{spill-remove} 高风险值；周围已经有多个值和 family 被染色；如果只做朴素逆序染色，
它可能直接落入 spill 或破坏已有 family 的颜色安排。真实 dump 的关键片段如下：

读这个例子前先区分几个区域。\texttt{coalesce} 是 family 构造结果。
\texttt{removal plan} 是 planner 的摘除顺序和 remove kind。
\texttt{select-trace} 是 reverse-select coloring 的逐值记录。
\texttt{alloc func} 是最终结果表。几个字段可以这样读：

\begin{itemize}
  \item \texttt{colors=3}：本轮可用颜色数，可以近似理解成 allocator 当前可分配的寄存器颜色预算；
  \item \texttt{candidates=2}：coalescing analysis 发现了两个可考虑合并/共色的 pair；
  \item \texttt{weight}：该 coalesce pair 的 affinity 权重，越高表示越值得尝试保留共色偏好；
  \item \texttt{can\_coalesce=yes}：当前 interference 关系允许这对值成为 coalescing 候选；
  \item \texttt{class ... root=... size=...}：family 代表值和 family 大小，例如 \texttt{ssa.2} 代表 \texttt{\{ssa.2, ssa.5\}}；
  \item \texttt{remove=simplify-remove}：该值在 planner 中按低压力 simplify 候选摘除；
  \item \texttt{remove=spill-remove}：该值在 planner 中按高风险 spill 候选摘除，但不等于最终一定 spill；
  \item \texttt{outcome=colored}：select 阶段普通染色成功；
  \item \texttt{outcome=optimistic-colored}：该值原本是 spill-remove 候选，但 select 阶段通过 optimistic 路线仍然找到了颜色；
  \item \texttt{color=N} / \texttt{slot=N}：最终落到某个颜色，或者真的进入某个 spill slot；
  \item \texttt{main-evict}：主 select 阶段为了当前值的颜色安排，选择某个 blocker 并尝试改色；本例中是把 \texttt{ssa.2} 挪到 \texttt{color1}。
\end{itemize}

在这个 witness 里，\texttt{ssa.0 priority=30 spill-intended} 的意思是：planner 曾经把
\texttt{ssa.0} 当作高风险 spill 候选摘掉，并记录了它的 spill priority；但最后
\texttt{outcome=optimistic-colored} 又说明 selector 成功把它救回了颜色。因此这里真正
值得看的不是单独的分数，而是
\texttt{remove}、\texttt{priority}、\texttt{outcome}、\texttt{main-evict} 这几个字段合在一起
描述的一次决策链。

\begin{lstlisting}[language={}]
coalesce manual_retry_blocker_family colors=3 candidates=2
  0: ssa.1 <-> ssa.3 weight=1 can_coalesce=yes
  1: ssa.2 <-> ssa.5 weight=3 can_coalesce=yes
  class ssa.1 -> root=ssa.1 size=2
  class ssa.2 -> root=ssa.2 size=2

select-trace colors=3 values=7
  0: ssa.4 remove=simplify-remove outcome=colored color=0
  1: ssa.2 remove=simplify-remove outcome=colored color=1
  2: ssa.5 remove=simplify-remove outcome=colored color=1
  3: ssa.6 remove=simplify-remove outcome=colored color=2
  4: ssa.1 remove=simplify-remove outcome=colored color=1
  5: ssa.3 remove=simplify-remove outcome=colored color=2
  6: ssa.0 remove=spill-remove outcome=optimistic-colored
     color=0 main-evict=ssa.2->color1

alloc func manual_retry_blocker_family colors=3 values=7
ssa.0 color=0 priority=30 spill-intended
ssa.1 color=1 priority=10
ssa.2 color=1 priority=10
ssa.3 color=2 priority=8
ssa.4 color=0 priority=5
ssa.5 color=1 priority=6
ssa.6 color=2 priority=40
\end{lstlisting}

这段 dump 的读法是：前半部分先说明 allocator 看见了两个 coalesced family；
后半部分按 reverse-select 顺序打印每个 value 的最终结果。最后一行最关键：
\texttt{ssa.0} 虽然是 \texttt{spill-remove} 摘下来的高风险值，但最后不是
confirmed spill，而是走到了 optimistic-colored。对应动作可以单独读成：

\begin{lstlisting}[language={}]
ssa.0 outcome=optimistic-colored
      main-evict=ssa.2->color1
\end{lstlisting}

也就是 selector 为了给 \texttt{ssa.0} 腾出可用安排，对 blocker \texttt{ssa.2} 执行了
recolor/evict 路线。

这里有一个细节需要讲清楚：select trace 中前面几行打印的是最终 allocation result，
不是“被 eviction 修改前”的历史快照。所以 \texttt{ssa.2} 那一行已经显示为最终
\texttt{color=1}；真正说明发生过 recolor 动作的是后面的 eviction metadata。
这里可以推导出 \texttt{ssa.2} 原先挡住的是 \texttt{color=0}：因为 \texttt{ssa.0}
最后 claim 到的颜色是 \texttt{color=0}，而 metadata 记录 blocker \texttt{ssa.2}
被挪到了 \texttt{color=1}。
这个 metadata 不是 dump 字符串里的猜测，而是 allocation result 里的结构化字段：

\begin{lstlisting}[language={}]
used_generic_eviction = yes
blocker_recolored = yes
blocker_new_color = 1
\end{lstlisting}

对应 regression 也会通过 query helper 断言：对 \texttt{ssa.0} 使用了 eviction，
blocker 是 \texttt{ssa.2}，并且 blocker 被 recolor 到 \texttt{color=1}。这个例子的价值
不在于它像源码程序一样自然，而在于它把 allocator 的策略决策锁成了可回归检查的 trace 字段。

和传统“放进栈里、逆序染色”的朴素版本相比，这里的差别不在于栈式 select 消失了，
而在于 select 时多了可解释的偏好和恢复动作。朴素实现通常只回答“这个点有没有颜色”；
这个 allocator 还能回答“为什么它想要这个颜色、是谁挡住了它、有没有把 blocker 挪开、
最后是不是真的 spill”。在热点代码里，这类策略可以减少无意义 move，避免把本来属于同一
value family 的值拆到不同位置；更重要的是，它让 allocator 的决策能被 dump 和 regression
test 直接锁住，而不是依赖主观判断。

这种设计带来的工程收益主要有：

\begin{itemize}
  \item move-related value 不再只是零散边，而是可以作为一组对象被观察；
  \item coalescing、freeze、spill 的原因可以通过 trace/report 解释；
  \item spill 后可以进行真实 rewrite，而不是只在报告里标记失败；
  \item no-move 函数可以走 fast path，避免为简单函数硬跑完整 family machinery；
  \item rewrite loop 可以只重分配被改写过的函数，降低不必要的编译时间。
\end{itemize}

\subsection{效果与边界}

\texttt{-perf} 模式在最终性能评测中通过了全部 20 个性能用例，正确性分数为
\texttt{100.00}。榜单记录为第 2 名，总用时 \texttt{103.72s}。单项结果如下：

\begin{table}[ht]
\centering
\footnotesize
\begin{tabular}{@{}lll@{}}
\toprule
用例 & 结果 & 用时 / s \\
\midrule
\texttt{00\_bitset1} & AC & 2.110806 \\
\texttt{01\_bitset2} & AC & 4.321101 \\
\texttt{02\_bitset3} & AC & 6.548924 \\
\texttt{03\_mm1} & AC & 7.327778 \\
\texttt{04\_mm2} & AC & 6.220852 \\
\texttt{05\_mm3} & AC & 4.664204 \\
\texttt{06\_mv1} & AC & 6.017828 \\
\texttt{07\_mv2} & AC & 3.883102 \\
\texttt{08\_mv3} & AC & 3.756907 \\
\texttt{09\_spmv1} & AC & 2.000937 \\
\texttt{10\_spmv2} & AC & 1.282645 \\
\texttt{11\_spmv3} & AC & 1.530294 \\
\texttt{12\_fft0} & AC & 2.949113 \\
\texttt{13\_fft1} & AC & 6.610910 \\
\texttt{14\_fft2} & AC & 6.063861 \\
\texttt{15\_transpose0} & AC & 2.295379 \\
\texttt{16\_transpose1} & AC & 2.728154 \\
\texttt{17\_transpose2} & AC & 5.738144 \\
\texttt{18\_brainfuck-bootstrap} & AC & 11.708512 \\
\texttt{19\_brainfuck-calculator} & AC & 15.961799 \\
\bottomrule
\end{tabular}
\caption{\texttt{-perf} 模式最终性能评测结果}
\end{table}

\texttt{-perf} 模式已经形成了比较完整的优化实验框架，但它并不宣称拥有一个
完全成熟的工业级优化器。当前更准确的定位是：

\begin{itemize}
  \item 已有不少真实、有回归保护的中端优化；
  \item 一些优化是 general pass，一些仍是 hotspot-specific pass；
  \item 部分 loop/induction 优化对具体 witness 有收益，但需要 A/B 判断是否保留；
  \item 编译时间同样是评价指标，不能为了小幅 runtime win 引入过高 compile-time cost。
\end{itemize}

\section{\texttt{-extension} 模式：语言扩展与 callable-object IR}

\texttt{-extension} 是项目中最有语言设计味道的一条线。它既是功能开关，也是
保守翻译管线开关。很多扩展特性只在该模式下开放，避免影响 \texttt{-riscv}
和 \texttt{-perf} 的课程基础目标。

\subsection{扩展特性概览}

当前 \texttt{-extension} 线覆盖了多类语言扩展：

\begin{itemize}
  \item control-flow extensions：\texttt{defer}、\texttt{fndefer}、\texttt{capdefer}；
  \item function-value family：non-capturing function values、local alias/reassignment、branch rebinding、function-valued parameters、returned callable、higher-order forwarding；
  \item closure/callable family：closure literal、local closure、returned closure、closure captures callable，以及后续的 callable-object IR；
  \item data/number extensions：pair、struct、aggregate 参数与返回，以及 conservative float slice。
\end{itemize}

这些特性不是一次性“全开”，而是按 source family 分阶段落地。每打开一个
新形状，就补 parser/semantic/IR/lower/runtime 等回归测试。项目里很多看似
特殊的小限制，本质上都是为了让每个 feature slice 都能稳定闭环。

\subsection{\texttt{defer} / \texttt{fndefer} / \texttt{capdefer}}

\texttt{defer} 这一组扩展主要解决“退出某个作用域时再执行一段代码”的问题。普通
\texttt{defer} 是 scope-exit 语义：注册点之后，当控制流离开当前作用域时，延迟语句按栈式顺序
执行。它和 Go 里的直觉类似，但项目里特意保守处理了控制流边界：defer body 中不能直接
\texttt{return}、\texttt{break} 或 \texttt{continue} 跳出外层；如果 \texttt{return expr}
遇到 active defer，会先固定返回值，再运行 deferred side effects。这样可以避免“返回值到底是
注册时、退出时还是副作用后”这类语义不确定性。

一个最小例子如下：

\begin{lstlisting}[language=C]
int main() {
  int x = 1;
  defer putint(x);
  x = 2;
  defer putint(3);
  return 0;
}
\end{lstlisting}

这里输出顺序是 \texttt{3, 2}：后注册的 \texttt{putint(3)} 先执行；前面的
\texttt{putint(x)} 读取的是退出时的 \texttt{x}，所以输出 2。

\texttt{fndefer} 是函数退出级别的 defer。它不绑定某个小 block 的退出，而是等函数即将返回时
统一运行；普通 block-level \texttt{defer} 会先按作用域 unwind，随后再运行函数级
\texttt{fndefer}。这个特性适合表达“无论从函数哪里返回，都要在最后做一件事”的模式，例如：

\begin{lstlisting}[language=C]
int g;
int main(int x) {
  int y = 0;
  fndefer { putint(g); putint(x); putint(y); }
  y = 3;
  return 0;
}
\end{lstlisting}

如果调用 \texttt{main(5)} 且全局 \texttt{g} 为 0，那么函数返回前会输出
\texttt{0, 5, 3}。这说明当前 \texttt{fndefer} 读的是函数退出时仍然稳定可见的名字，而不是
注册瞬间的快照。

当前 \texttt{fndefer} 不是任意闭包捕获系统，而是一个保守的函数退出注册机制。它要求引用的名字
在函数级视角下稳定可见；嵌套 block local、loop local、尚未声明的 later local 都会被拒绝。
\texttt{fndefer} body 中也不允许再嵌套 \texttt{defer}/\texttt{fndefer} 或制造新的
\texttt{return}/\texttt{break}/\texttt{continue} 控制转移。和循环有关的限制更微妙：
项目后来支持了部分动态注册形状，例如在 \texttt{if} 中按路径注册一次 \texttt{fndefer}，
也支持一个 loop-nested 注册点被重复执行；实现上会用隐藏计数器记录该 site 执行了多少次，
并在函数退出时 replay 对应次数。但当前仍只支持每个函数一个这样的 loop-nested site，
不是完整的运行时 defer stack。

例如，下面这种条件注册是可以表达的：

\begin{lstlisting}[language=C]
int main() {
  int x = 1;
  if (x) { fndefer putint(x); }
  x = 4;
  return 0;
}
\end{lstlisting}

输出是 \texttt{4}。但下面这种引用 loop-local 的写法会被拒绝：

\begin{lstlisting}[language=C]
int main() {
  for (int i = 0; i < 3; i = i + 1) {
    fndefer putint(i); // rejected: loop-local name
  }
  return 0;
}
\end{lstlisting}

还有一种拒绝和名字作用域无关，而是来自动态注册机制本身：同一个函数里出现多个
loop-nested \texttt{fndefer} site 时，当前实现会拒绝，即使它们只引用全局变量：

\begin{lstlisting}[language=C]
int a[8];
int main() {
  for (int i = 0; i < 8; i = i + 1) {
    fndefer putint(a[2]);
    fndefer putint(a[3]); // rejected: second loop site
  }
  return 0;
}
\end{lstlisting}

对应的设计取舍是：单个重复注册点可以被编译成“计数器 + 退出时循环 replay”，而多个动态
注册点需要更接近 runtime stack 的顺序记录。这里先 gate 住，避免半支持导致退出顺序不可靠。

\texttt{capdefer} 则是对这个缺口的一个小而明确的补充：它显式把表达式按值捕获到新名字里，
再在退出时执行延迟语句。例如：

\begin{lstlisting}[language=C]
int main() {
  int x = 1;
  capdefer(y = x) putint(y);
  x = 2;
  return 0;
}
\end{lstlisting}

这里延迟执行时打印的是注册时捕获的 \texttt{y}，而不是退出时已经变成 2 的 \texttt{x}。
因此，\texttt{defer} 偏 exit-time state，\texttt{capdefer} 偏 registration-time snapshot。
上面例子的输出是 \texttt{1}；如果换成普通 \texttt{defer putint(x)}，输出就会是
\texttt{2}。\texttt{capdefer} 同样保持保守边界：它只在 \texttt{-extension} 下开放，循环注册和
引用嵌套 block-local 的危险形状会被拒绝。例如：

\begin{lstlisting}[language=C]
int main() {
  int x = 1;
  for (; x < 2; x = x + 1) {
    capdefer(y = x) putint(y); // rejected: dynamic capdefer registration
  }
  return 0;
}
\end{lstlisting}

\subsection{\texttt{pair}、\texttt{struct} 与 aggregate}

\texttt{pair}/\texttt{struct} 这一组扩展表面上是在给语言增加复合值，工程上更重要的点则是：
项目在不引入完整对象 ABI 的前提下，先让复合值以保守、可验证的方式流动起来。

\texttt{pair} 是最小的 builtin aggregate，支持局部声明、\texttt{\{1, 2\}} 初始化、
copy-init/assignment，以及 \texttt{.first}/\texttt{.second} 字段访问。\texttt{struct}
则把这条线推进到 named aggregate：用户可以在顶层声明类型，再在局部创建对象、初始化、
赋值和访问字段。例如：

\begin{lstlisting}[language=C]
struct Point { int x; int y; int z; };

int main() {
  pair p = {1, 2};
  struct Point q = {3, 4, 5};
  q.x = p.first + q.y;
  return q.x + p.second;
}
\end{lstlisting}

这里的实现并不是把 \texttt{p} 或 \texttt{q} 变成一个运行时对象，而是分层识别、再递归展平成
hidden scalar slots。具体可以拆成四步。

第一步，parser/AST 保留两类信息：顶层 \texttt{struct} 定义会进入 \texttt{AstProgram} 的
struct type table；字段访问则是普通的 member expression，形如 \texttt{base.field}。所以
\texttt{b.p.second} 在 AST 中不是一个扁平字符串，而是一条 member chain。

第二步，semantic 根据 scope 中记录的 declaration kind 和 type name 识别 aggregate root。
\texttt{pair} 的合法字段固定是 \texttt{first}/\texttt{second}；\texttt{struct} 则按
struct type table 查字段名。如果字段本身还是 \texttt{pair} 或另一个 known \texttt{struct}，
semantic 返回的仍然是 aggregate type descriptor，于是下一层 \texttt{.second} 或 \texttt{.x}
还能继续递归解析。

第三步，IR lowering 给每个 aggregate local 分配一段连续 scalar locals，并在 lowering scope
里记录：

\begin{lstlisting}[language={}]
aggregate binding =
  first_local_id + aggregate_slot_count + aggregate_kind + type_name
\end{lstlisting}

因此 lowering 心智模型更接近：

\begin{lstlisting}[language={}]
pair p = {1, 2}
  -> p.first_slot  = 1
  -> p.second_slot = 2

struct Point q = {3, 4, 5}
  -> q.x_slot = 3
  -> q.y_slot = 4
  -> q.z_slot = 5
\end{lstlisting}

第四步，字段访问才真正“用起来”：lowering 先解析 member chain 的 base 是哪个 aggregate
binding，再根据 struct 字段顺序计算 flattened offset。对 \texttt{pair} 来说，
\texttt{first} 是 offset 0，\texttt{second} 是 offset 1；对 \texttt{struct} 来说，某个字段的
offset 等于它前面所有字段 flatten 后的 slot 数之和。

这个策略自然支持嵌套 aggregate，例如：

\begin{lstlisting}[language=C]
struct Box { pair p; int tag; };

int main() {
  struct Box b = {{1, 2}, 7};
  pair q = b.p;
  return q.second + b.tag;
}
\end{lstlisting}

在这个例子里，\texttt{Box} 的 flatten layout 可以理解为：

\begin{lstlisting}[language={}]
b.p.first  -> slot 0
b.p.second -> slot 1
b.tag      -> slot 2
\end{lstlisting}

所以 \texttt{b.p.second} 的解析过程是：先确认 \texttt{b} 是 \texttt{struct Box}，再确认字段
\texttt{p} 是 \texttt{pair}，于是 \texttt{b.p} 变成一个 aggregate-valued sub-binding；
最后 \texttt{second} 映射到这个 sub-binding 的第 1 个 scalar leaf。copy-init 和 assignment
也是逐 leaf 搬运，而不是复制一个 opaque object。真实 IR witness 里也能看到这种形状，例如
\texttt{struct Mid \{ pair p; int z; \}} 会展开出类似 \texttt{m\$p\$0}、\texttt{m\$p\$1}、
\texttt{m\$z} 的 locals。

这条线后来还保守地打开了函数边界：\texttt{pair} 和 known \texttt{struct Name} 可以作为
参数和返回值出现在若干受控 family 中，例如：

\begin{lstlisting}[language=C]
pair id(pair p) { return p; }
int sum(pair p) { return p.first + p.second; }
pair mk() {
  pair p = {10, 20};
  return p;
}
int main() {
  pair q = id(mk());
  return sum(q);
}
\end{lstlisting}

这仍然不是完整 C aggregate ABI。更准确地说，caller/callee 共享同一套 flatten contract：
aggregate 参数会展开成 hidden scalar parameter sequence，aggregate 返回则通过 hidden result
slots 写回。也就是说，跨函数边界流动的是 leaf slots，不是一个任意位置都可作为普通表达式消费的
first-class aggregate object。

因此当前边界也很明确：aggregate globals、aggregate arrays、任意 scalar context 中的 plain
aggregate value，以及直接对调用结果取字段的形状仍然会拒绝。比如 \texttt{mk().second} 当前不会
被当作“临时 aggregate 对象上读字段”自动打开；更推荐、也更符合当前 lowering contract 的写法是：

\begin{lstlisting}[language=C]
pair q = mk();
return q.second;
\end{lstlisting}

这正是本项目处理语言扩展的一贯风格：先把 type descriptor、field lookup、flatten order、
copy/return/param compatibility 这些基础设施收稳，再决定是否进入真正的一等 aggregate ABI。
这种小步推进可以让扩展能力保持可验证，而不是在 ABI 尚未稳定时扩大语义表面。

\subsection{函数值与闭包}

函数值主线是 \texttt{-extension} 最值得重点介绍、也最容易误解的部分。它的目标不是一次性实现
完整 C 函数指针或 ML 风格闭包运行时，而是在当前编译器架构里逐步回答三个问题：

\begin{enumerate}[label=\arabic*.]
  \item 函数是否可以作为值传给另一个函数；
  \item callable 离开定义点之后，在返回、赋值、转发、捕获时是否还能被正确恢复；
  \item 这些恢复动作是否能从 lowering context 里的隐式约定，收敛成 IR/verifier 可见的对象合约。
\end{enumerate}

因此实现过程不是先引入完全通用的 function object，而是分阶段推进。

\paragraph{第一阶段：non-capturing specialization。}
最早支持的是无捕获、目标静态已知的函数值参数：

\begin{lstlisting}[language=C]
int add1(int x) { return x + 1; }
int apply(int f(int), int x) { return f(x); }
int main() { return apply(add1, 41); }
\end{lstlisting}

这一层的核心不是“已经有通用函数指针 ABI”，而是 specialization：调用点知道
\texttt{f=add1}，于是构造一个 specialized helper。心智模型大致是：

\begin{lstlisting}[language={}]
apply(add1, 41)
  -> apply__fv_0_add1(41)
  -> call add1(41)
\end{lstlisting}

后续的 local alias、forwarding、多 function-valued parameters、zero-argument、\texttt{void}
return 等 sibling 都是在这条线上扩张。比如：

\begin{lstlisting}[language=C]
int compose(int f(int), int g(int), int x) {
  return f(g(x));
}
\end{lstlisting}

如果 \texttt{f} 和 \texttt{g} 在调用点都是已知 top-level 函数，lowering 就可以生成对应组合的
specialized helper。

\paragraph{第二阶段：dynamic noncapturing 与 returned callable。}
一旦出现运行时选择，单纯静态 specialization 就不够了：

\begin{lstlisting}[language=C]
int add2(int x) { return x + 2; }
int main(int c) {
  int f(int) = add1;
  if (c) f = add2;
  return f(40);
}
\end{lstlisting}

这类值当前用 target family/tag 表达。编译期先为这个 local 推出一个有限 target set，例如
\texttt{\{add1, add2\}}，并给每个 target 一个内部 tag；local 本身会多出一个 hidden tag slot，
dump 里常见名字类似 \texttt{f\$ftag}。初始化和赋值并不复制一个真实函数对象，而是更新这个
tag：

\begin{lstlisting}[language={}]
int f(int) = add1;  -> f$ftag = tag(add1)
f = add2;           -> f$ftag = tag(add2)
\end{lstlisting}

调用 \texttt{f(40)} 时，lowering 读取 \texttt{f\$ftag}，按 tag 走不同分支。每个分支仍然落到
编译期已经生成好的 specialized leaf：

\begin{lstlisting}[language={}]
if f$ftag == tag(add1):
    call __fnwrap_add1(..., 40)
else if f$ftag == tag(add2):
    call __fnwrap_add2(..., 40)
\end{lstlisting}

如果这个动态 local 被传给 \texttt{apply(f, 40)}，也是同一个思路：先看 \texttt{f\$ftag}，
再选择某个已生成的 \texttt{apply} specialized helper：

\begin{lstlisting}[language={}]
tag(add1) -> apply__fv_0_add1(40)
tag(add2) -> apply__fv_0_add2(40)
\end{lstlisting}

所以第二阶段仍然不是 generic runtime function pointer；它是“有限 target family + hidden tag
+ branch-selected specialized helper”。

再往前一步，函数值还可以跨 return boundary：

\begin{lstlisting}[language=C]
int pick(int c)(int) {
  if (c) return add2;
  return add1;
}
int main() {
  int f(int) = pick(1);
  return f(40);
}
\end{lstlisting}

这里真正的工作不是 parse \texttt{int pick()(int)}，而是 producer 返回一份 conservative
callable payload，caller 再 materialize 本地 callable view。相关修复集中在这里：
同一个 returned producer 不能重复求值，ternary/passthrough/wrapper 里要找回真正 producer，
多个 function-valued actual argument 还要组合它们的 dispatch。

\paragraph{第三阶段：closure 与 capture payload。}
closure 让问题从“返回哪个函数”升级为“返回哪个函数以及哪些捕获值”。第二阶段里
\texttt{f\$ftag} 基本只需要回答“现在是哪一个 non-capturing target”；第三阶段则要同时回答：
“现在是哪一个 closure helper”，以及“这个 helper 运行时要看见哪一组 captured values”。本地
closure 的基本形状是：

\begin{lstlisting}[language=C]
int main() {
  int x = 3;
  int f(int) = closure [x] int (int y) { return x + y; };
  return f(4);
}
\end{lstlisting}

lowering 会先把 closure literal 拆成一个 helper/wrapper，再把 capture list 拆成 hidden
payload。对上面的例子，可以把 local \texttt{f} 想成被展开成了两类隐藏状态：

\begin{lstlisting}[language={}]
f$ftag          = tag(main__closure_f)
f$closurecap$0  = x
\end{lstlisting}

\texttt{f\$ftag} 表示这个 local 当前绑定到哪一个 closure family；\texttt{f\$closurecap\$0}
保存第 0 个捕获值，也就是初始化时的 \texttt{x}。这仍然是 by-value capture：之后外层
\texttt{x} 如果再被赋值，不会自动改写已经保存进 payload 的值。

调用 \texttt{f(4)} 时，lowering 不会再去“按名字查外层的 \texttt{x}”。它先根据
\texttt{f\$ftag} 找到 closure wrapper，再把 capture slot 当作 hidden env 传进去：

\begin{lstlisting}[language={}]
env = f$closurecap$0
obj = fn_make __fnwrap_closure_main__closure_f, env, shape(int)->int
call_indirect obj(4)
\end{lstlisting}

wrapper 的物理调用形式可以理解成 \texttt{code(env, visible\_args...)}，所以 closure body 里的
\texttt{x} 是从 env/capture payload 中恢复出来的。这里的关键变化是：source 里只有
\texttt{f(4)} 一个参数，但 wrapper 实际上会多吃一个 hidden env；因此后续 verifier 和 lower
阶段必须一直区分 hidden 参数和 source-visible 参数。

returned closure 进一步要求 payload 跨函数边界：

\begin{lstlisting}[language=C]
int make(int x)(int) {
  return closure [x] int (int y) { return x + y; };
}
int main() {
  int f(int) = make(3);
  return f(4);
}
\end{lstlisting}

这个例子里，\texttt{make(3)} 不能只返回一个 tag。它至少要把 closure family identity 和捕获值
\texttt{x=3} 一起带回 caller。当前实现仍然避免完整 heap closure runtime，而是使用保守的
hidden return payload。可以将其概括为：

\begin{lstlisting}[language={}]
callee make:
  __retclosure_tag       = tag(make__retclosure_...)
  __retclosure_cap_0     = x

caller main:
  f$ftag                 = __retclosure_tag
  f$closurecap$0         = __retclosure_cap_0
\end{lstlisting}

真实 dump 中的名字会更工程化，常见前缀包括：

\begin{lstlisting}[language={}]
__retclosure_*
__retclosure_argcap_*
__retclosure_immediate_*
__retclosure_specarg_*
\end{lstlisting}

这些名字的职责是把 returned producer 产生的 family tag 和 capture slots 稳定地交给
caller。之后 \texttt{f(4)} 就和本地 closure 调用走
同一条 wrapper / callable-object 路径。

当前还支持一些更复杂但仍保守的 family，例如 closure 捕获 callable：

\begin{lstlisting}[language=C]
int wrap(int f(int), int x) {
  int g(int) = closure [f] int (int y) { return f(y); };
  return g(x);
}
\end{lstlisting}

这时 capture payload 里不只有 scalar，还可能有 function-value tag / callable payload。也就是说，
\texttt{closure [f]} 捕获的不是一个普通整数值，而是“当前 \texttt{f} 代表哪个 callable family”
的内部 payload；closure helper 进入 body 后还要恢复 \texttt{f(y)} 这个调用。实现上因此需要：

\begin{lstlisting}[language={}]
capture time:
  g$closurecap$0 = f$ftag or rebuilt callable payload

inside closure helper:
  recover captured callable
  dispatch captured f(y)
\end{lstlisting}

这也是为什么这条线后来不再适合只依赖 helper 名称拼接：payload 可能来自 local binding、
returned producer、passthrough wrapper 或多个 function-valued actual argument，必须有统一的
materialize/rebuild 规则，否则容易出现同一个 producer 被重复求值、不同 payload prefix 指向
不同临时槽等真实错误。

\paragraph{第四阶段：callable-object IR。}
随着 returned callable、returned closure、closure captures callable、higher-order passthrough
越来越多，项目把共同抽象收敛成 canonical IR 中显式可见的 callable object。前三阶段仍然可以
理解成在 source lowering 里维护 hidden tag、capture slot 和 helper；第四阶段要做的是把这些
约定变成 IR 里能 dump、能 verify、能继续 lowering 的对象：

\[
\text{callable object} = (\text{code}, \text{env}, \text{shape})
\]

它对应三个字段：\texttt{code} 是 wrapper/helper 入口，\texttt{env} 是 hidden capture payload，
\texttt{shape} 是 source-visible signature。non-capturing function、local closure 和 returned
closure 都可以被映射到同一种构造动作：

\begin{lstlisting}[language={}]
static function add1:
  code = __fnwrap_add1
  env  = 0

local / returned closure:
  code = __fnwrap_closure_...
  env  = f$closurecap$0 or returned capture payload

object:
  obj = fn_make code, env, shape
\end{lstlisting}

在 canonical IR 中，最小调用形态是：

\begin{lstlisting}[language={}]
shape shape.0(int) -> int
tmp.0 = fn_make __fnwrap_add1, 0, shape.0
tmp.1 = call_indirect tmp.0(41)
\end{lstlisting}

这里 \texttt{shape} 专门描述 visible signature，例如 \texttt{(int)->int}；它不包含 hidden env。
同时，indirect call 也只接收 visible args。这个区别很重要：closure 可能有 env，但调用参数数量
仍然只看 source-level 形参。因此 verifier 可以检查三件事：

\begin{itemize}
  \item \texttt{fn\_make} 指向的 wrapper/helper 是否存在；
  \item wrapper 的物理参数是否等于 \texttt{hidden env + visible args}；
  \item \texttt{call\_indirect} 的 visible arg count 是否匹配 \texttt{shape}。
\end{itemize}

这也是 zero-arg closure、void closure 等 case 容易出 bug 的原因：env/capture slot 数量可以
不为 0，但 visible arg count 必须仍然是 0。

lower-IR/backend 不一定一直保留 \texttt{call\_indirect} 文本形状。如果三元组已经明确，可以桥接成：

\begin{lstlisting}[language={}]
call __fnwrap_closure_...(env, visible_args...)
\end{lstlisting}

所以正文或 dump 里看到 direct wrapper call，并不代表第四阶段没有发生；它通常表示 canonical IR
先建立了显式 callable-object contract，后续层再把它安全地折回普通 call。

这条 IR 线的意义在于：callable 不再只藏在 lowering context 和 payload slot 约定里，而是变成
dump 可展示、verifier 可检查、lower/backend 可桥接的显式对象。后端如果证明
\texttt{code/env} 已知，可以再把它折回直接 wrapper call；这不是倒退，而是“IR 合约显式化，
后层允许优化”的正常分工。

\paragraph{higher-order：有限层显式签名，而不是任意泛型。}
项目还支持了一批 higher-order 形状，例如函数参数本身也接收函数参数：

\begin{lstlisting}[language=C]
int apply(int f(int), int x){ return f(x); }
int pass(int h(int f(int), int x), int f(int), int x){
  return h(f, x);
}
\end{lstlisting}

这一类依赖 parser/semantic 保留 nested function signature metadata；semantic 递归比较
return kind、parameter count 和 nested function parameter shape；lowering 再用 passthrough
分析和 recursive dispatch，把 outer callable binding 传到 leaf callable-object dispatch。
所以它不是自动泛型 lambda，而是“有限层显式 shape + specialization/callable-object collapse”。

\paragraph{当前能写什么。}
概括地说，当前 \texttt{-extension} 已能覆盖这些代表性 family：

\begin{itemize}
  \item top-level/builtin function 作为 function-valued parameter，含多参数、zero-arg、\texttt{void} sibling；
  \item local function-typed alias、部分 reassignment、branch rebinding 和 dynamic tag dispatch；
  \item noncapturing function-valued return，以及 bind-and-call / immediate-call 消费；
  \item local closure、returned closure、multi-capture closure，以及部分 closure reassignment / alias；
  \item closure literal 或 returned closure 作为 function-valued actual argument；
  \item closure 捕获静态或参数传入的 function value；
  \item 显式 nested signature 的 second/third/... order passthrough family；
  \item canonical IR 上的 \texttt{shape + fn\_make + call\_indirect} callable-object 路径。
\end{itemize}

\paragraph{当前不能写什么。}
同样重要的是边界。当前项目仍不声称支持：

\begin{itemize}
  \item 完整 generic function pointer ABI 或所有 callable 都可任意存取的 runtime；
  \item callable value 存进 global、array、aggregate field 后再自然流动；
  \item 任意深度、任意表达式形态的 CPS/continuation 程序；
  \item polymorphic/auto lambda、自递归 callable、真正无限层 higher-order type inference；
  \item 所有 closure literal + function-valued-parameter 组合，尤其是需要完整动态 continuation stack 的形状；
  \item 完整 heap escaping closure lifetime 模型。
\end{itemize}

换句话说，当前实现拥有的是一条很宽的 conservative function-value mainline，而不是最终形态的
first-class function runtime。这个边界反而是设计的一部分：支持的 family 用 regression 锁住；
暂时不支持的 family 尽量在 semantic gate 处明确拒绝，避免把未完成能力暴露成不稳定行为。

\section{端到端翻译案例}

为了避免报告只停留在结构描述上，本节选取几段小而有代表性的
\texttt{-extension} 程序，展示它们在不同阶段中的关键形状。这里不把所有扩展硬塞进同一个
综合样例，因为某些组合本来就是当前 semantic gate 会保守拒绝的；更清楚的方式是：
用一个主 showcase 展示 aggregate + closure + callable-object 主线，再用一个补充 showcase 展示
defer/capdefer/fndefer 的控制流改写。正文只展示关键片段，以避免大段 dump 干扰主线叙述。

\subsection{主 showcase：aggregate + returned closure}

\begin{lstlisting}[language=C]
struct Box { pair p; int bias; };

int sum_box(struct Box b) {
  return b.p.first + b.p.second + b.bias;
}
struct Box mkbox(int a, int b, int c) {
  struct Box box = {{a, b}, c};
  return box;
}

int apply(int f(int), int x) { return f(x); }
int pass(int h(int f(int), int x), int f(int), int x) {
  return h(f, x);
}
int make(int base)(int) {
  return closure [base] int (int y) { return base + y; };
}
int main() {
  struct Box b = mkbox(1, 2, 3);
  int f(int) = make(sum_box(b));
  return pass(apply, f, 4);
}
\end{lstlisting}

这个例子刻意覆盖几类 \texttt{-extension} 形状，但仍保持可读。\texttt{Box} 里嵌了
\texttt{pair}；\texttt{mkbox} 通过 hidden return slots 返回 aggregate；\texttt{sum\_box}
通过 flattened aggregate 参数读取字段。随后程序用 \texttt{sum\_box} 的结果构造 returned
closure，再通过 \texttt{pass}/\texttt{apply} 展示 higher-order passthrough 被折到 leaf
callable dispatch。

\subsection{aggregate lowering 关键片段}

\begin{lstlisting}[language={}]
func sum_box(b.0, b.1, b.2) {
  bb.0:
    tmp.0 = mov b.0
    tmp.1 = mov b.1
    tmp.2 = add tmp.0, tmp.1
    tmp.3 = mov b.2
    tmp.4 = add tmp.2, tmp.3
    ret tmp.4
}

func mkbox(__ret0.0, __ret1.1, __ret2.2, a.3, b.4, c.5) {
  bb.0:
    box$p$0.6 = mov a.3
    box$p$1.7 = mov b.4
    box$bias.8 = mov c.5
    0 = store_indirect __ret0.0, box$p$0.6
    0 = store_indirect __ret1.1, box$p$1.7
    0 = store_indirect __ret2.2, box$bias.8
    ret
}
\end{lstlisting}

这段说明 aggregate 当前不是走完整 C ABI 的“结构体整体对象”，而是被拆成 leaf scalar slots。
\texttt{struct Box \{ pair p; int bias; \}} 被展开为 \texttt{p.first}、\texttt{p.second}、
\texttt{bias} 三个槽；作为参数时表现为 \texttt{b.0/b.1/b.2}，作为返回值时则通过
\texttt{\_\_ret0/\_\_ret1/\_\_ret2} 这类 hidden return slots 写回 caller。

\subsection{callable-object 关键片段}

\begin{lstlisting}[language={}]
shape shape.0(int) -> int

func main() {
  bb.0:
    tmp.0 = addr_local local.0
    tmp.1 = addr_local local.1
    tmp.2 = addr_local local.2
    tmp.3 = call mkbox(tmp.0, tmp.1, tmp.2, 1, 2, 3)
    tmp.4 = addr_local local.3
    tmp.5 = call sum_box(b$p$0.0, b$p$1.1, b$bias.2)
    tmp.6 = call make(tmp.4, tmp.5)
    f$ftag.4 = mov 1
    __closure_env_0.5 = mov f$closurecap$0.3
    tmp.7 = addr_local local.5
    tmp.8 = fn_make __fnwrap_closure_make__retclosure_7_33_1,
                    tmp.7, shape.0
    tmp.9 = call_indirect tmp.8(4)
    ret tmp.9
}
\end{lstlisting}

这段把几条线连在一起了：\texttt{mkbox} 先把 aggregate 写到 caller 准备的 hidden slots；
\texttt{sum\_box(b)} 再按 flattened slots 读出字段；\texttt{make(...)} 把求出的
\texttt{base} 保存为 returned closure payload；最后 \texttt{fn\_make} 与
\texttt{call\_indirect} 把 closure wrapper、环境地址和 callable shape 合成 verifier 可见的
callable object。这里源码写的是 \texttt{pass(apply, f, 4)}，但 passthrough 分析已经把外层
higher-order shell 折到 leaf callable dispatch，所以 IR 里直接出现最终 closure 调用路径。

\subsection{lower IR 与 RISC-V 关键片段}

\begin{lstlisting}[language={}]
func main() {
  bb.0:
    tmp.0 = addr_local b$p$0.0
    tmp.1 = addr_local b$p$1.1
    tmp.2 = addr_local b$bias.2
    call mkbox(tmp.0, tmp.1, tmp.2, 1, 2, 3)
    tmp.11 = load_local b$p$0.0
    tmp.12 = load_local b$p$1.1
    tmp.13 = load_local b$bias.2
    tmp.5 = call sum_box(tmp.11, tmp.12, tmp.13)
    call make(addr_local f$closurecap$0.3, tmp.5)
    store_local f$ftag.4, 1
    tmp.9 = call __fnwrap_closure_make__retclosure_7_33_1(
        addr_local __closure_env_0.5, 4)
    ret tmp.9
}
\end{lstlisting}

lower IR 的第一件事，是把 local/global 访问显式拆成
\texttt{load\_*} 与 \texttt{store\_*}。与此同时，当前 slice 会把
\texttt{call\_indirect} 桥接为 wrapper call。最终输出的 RISC-V 片段也能看到
同样的结构：

\begin{lstlisting}[language={}]
main:
  mv a0, sp
  addi a2, sp, 4
  addi a1, sp, 8
  call mkbox
  call sum_box
  mv a1, a0
  addi a0, sp, 12
  call make
  ...
  li a1, 4
  call __fnwrap_closure_make__retclosure_7_33_1
  ret
\end{lstlisting}

这个案例体现了项目的一个核心风格：源码、IR、lower IR 和汇编之间不是黑箱跳跃，
而是每一层都有能解释的中间形状。对复杂语言特性来说，这比“最后能跑”更重要；
最终运行结果仍然必要，但中间层契约也必须能够解释。

\subsection{defer / capdefer / fndefer 展示}

另一类扩展不主要改变值的表示，而是改变 return 前的控制流。一个小例子是：

\begin{lstlisting}[language=C]
int main() {
  int x = 1;
  defer { x = x + 10; }
  capdefer(y = x) putint(y);
  x = 2;
  fndefer putint(x);
  return x;
}
\end{lstlisting}

对应 canonical IR 的关键形状如下：

\begin{lstlisting}[language={}]
func main() {
  bb.0:
    __fndefer_count_0.0 = mov 0
    x.1 = mov 1
    __capdefer_0_y.2 = mov x.1
    x.1 = mov 2
    __fndefer_count_0.0 = mov __fndefer_count_0.0 + 1
    __fndefer_ret.3 = mov x.1
    x.1 = mov x.1 + 10
    call putint(__capdefer_0_y.2)
    __fndefer_i.4 = mov 0
    jmp replay_fndefer
  replay_fndefer:
    if __fndefer_i.4 < __fndefer_count_0.0:
      call putint(x.1)
      __fndefer_i.4 = __fndefer_i.4 + 1
      jmp replay_fndefer
    ret __fndefer_ret.3
}
\end{lstlisting}

这段展示了三者的差别：\texttt{capdefer} 在注册点保存 snapshot，所以打印的是旧的
\texttt{x=1}；普通 \texttt{defer} 在 return 前展开并修改当前 \texttt{x}；\texttt{fndefer}
则通过 hidden counter 和 replay loop 表示函数退出级注册。注意返回值先被保存到
\texttt{\_\_fndefer\_ret}，所以 defer 修改 \texttt{x} 不会意外改变本次 \texttt{return x}
已经确定的返回值。

\section{严格返回路径审计}

除了三个主模式外，项目还提供了一个可叠加的 strict gate：

\begin{lstlisting}[language={}]
--enforce-all-paths-return-check
\end{lstlisting}

它不是第四种编译模式，而是一个更严格的语义审计开关。默认 driver 会跳过这个检查，是为了兼容
课程测试中较宽松的非 \texttt{void} 函数写法；但在需要审计控制流时，可以显式打开它，让
semantic 拒绝没有 all-path return 证明的函数。

\subsection{为什么需要单独做这个检查}

只检查“函数里出现过 \texttt{return}”是不够的。下面这个函数有 \texttt{return}，但显然不是
所有路径都返回：

\begin{lstlisting}[language=C]
int f(int x) {
  if (x) return 1;
}
\end{lstlisting}

当 \texttt{x==0} 时，程序会直接落到函数末尾。对非 \texttt{void} 函数来说，这就是未定义的
返回值来源。课程基础模式为了兼容性默认放宽；strict gate 则把它作为明确的语义错误报告为
\texttt{SEMA-CF-001}。

\subsection{实现思路}

这个检查不是运行时检查，也不是完整停机证明，而是在语义分析阶段对 statement AST 做控制流摘要。
实现上会计算两个核心信息：

\begin{lstlisting}[language={}]
returns_all_paths = guaranteed_return && !may_fallthrough
\end{lstlisting}

\texttt{guaranteed\_return} 表示当前语句树是否形成必然返回边界；
\texttt{may\_fallthrough} 表示是否仍有路径能走到函数末尾而不返回。只有非
\texttt{void} 函数满足上式，strict 模式才接受。这个摘要会递归处理 compound、if/else、
while、for、break/continue 等控制流；它的目标不是证明所有程序性质，而是保守、稳定、
可解释。

\subsection{循环与不可达 return}

循环形状能较清楚体现这个检查的价值。例如：

\begin{lstlisting}[language=C]
int f(int a) {
  while (1) {
    if (a) return 1;
  }
}
\end{lstlisting}

它后面即使补一个不可达的 \texttt{return}，也不能伪装成所有路径都返回；因为
\texttt{a==0} 时会进入一个静态可见的非返回路径。相对地，如果循环体内部已经覆盖所有分支：

\begin{lstlisting}[language=C]
int f(int a) {
  while (1) {
    if (a) return 1;
    else return 2;
  }
}
\end{lstlisting}

则 strict gate 会接受。类似地，\texttt{for(;;)} 中如果所有可达分支都会返回，也可以通过；
如果只有部分分支返回，则会被拒绝。这一点避免了用无限循环后方的不可达
\texttt{return} 掩盖真实的非返回路径。

\subsection{边界与测试价值}

这个检查拒绝的不是所有潜在风险程序，而是在当前保守分析下不能证明 all-path return
的程序。它也不会声称解决停机问题：值依赖、调用依赖、复杂循环退出等情况会保持谨慎，避免把
合法程序误杀。项目的测试中也专门区分了稳定 reject、保守 accept 和 deterministic loop reject
几类 case，防止后续语义分析改动把这个 gate 变得过宽或过窄。

从工程角度看，all-paths-return 的意义不只是多报一个错。它给前端控制流分析提供了一个小而清楚
的证明任务：哪些路径会返回，哪些路径可能 fall through，哪些循环不能被不可达的后续
\texttt{return} 误判为已覆盖。编译器可以允许默认模式保持兼容性，但 strict 模式必须提供更严格的
语义审计。

\section{Vibe 工作流与质量保障}

本项目开发过程中采用了比较典型的快速迭代式 vibe coding 工作流：先围绕一个具体
witness 或 feature slice 形成假设，再实现最小改动，随后立即通过测试和实际运行验证。
这种方式的好处是推进快，风险是容易在复杂 lowering 中产生仅由最终现象支撑的错误判断。
因此项目把质量保障放在非常靠前的位置。

\subsection{AI agent 协作设定}

本项目的一个特殊之处，是它几乎完全通过 AI agent 协作完成代码层面的推进。
人的角色并不是逐行写 C 代码，而是负责路线、目标、优先级和大体实现方法：
决定当前要攻哪一簇问题、什么算真正修好、哪些风险需要暂停，以及什么时候必须跑全量测试。
在实际开发中，项目负责人参与了架构决策、验收和方向控制，但基本不直接书写实现代码；
几十万行级别的实现主要由 agent 根据目标拆解、阅读代码并落地。

协作大致分成三类 agent：

\begin{itemize}
  \item 实现 agent：阅读当前权威文档和代码，负责在 \texttt{src/}、\texttt{include/}、\texttt{tests/}、\texttt{docs/} 内实现功能、修 bug、补回归测试；
  \item review/lesson agent：阅读代码和测试结果，做代码 review、归纳项目结构，并将复杂实现压缩成 \texttt{lesson/} 中更容易复用的说明；
  \item planning/hacking agent：阅读代码、lesson 和路线文档，精炼下一步计划，同时主动寻找当前实现或规划中可能有问题的地方。
\end{itemize}

这些 agent 不是随意开工。仓库中的 \texttt{AGENTS.md} 定义了启动顺序、默认角色、
作用域边界、测试纪律和版本规则。例如 agent 开始工作前需要读取
\texttt{AGENTS.md}、\texttt{docs/ENGINEERING\_MEMORY.md} 和
\texttt{docs/NEXT\_STEPS.md}；如果进入优化线，还要读取优化计划和 reusable pass
待办。它还明确规定了不能随意修改 \texttt{lesson/}、不能提交已知 broken code、
性能改动必须做 A/B、长命令不能随便打断等规则。

模型侧主要使用 GPT 5.4/5.5 级别的 coding/reasoning agent 进行 vibe coding。
这里的 ``vibe'' 并不等于没有纪律，恰恰相反：人的路线判断负责“往哪里走”，agent
负责“怎么把代码走出来”，而 verifier、regression、runtime test 和 full suite
负责判断改动是否可接受。与传统的手写代码为主、工具辅助检查不同，本项目更强调
目标设定、实现自动化和验证闭环之间的分工：agent 可以快速生成和修改实现，但最终状态
必须由可重复的测试与诊断结果确认。

\subsection{测试矩阵}

项目测试不是单一入口，而是覆盖多个层次：

\begin{itemize}
  \item lexer/parser/semantic regression：锁前端语法和语义边界；
  \item IR/lower IR regression：锁 lowering 形状和 verifier contract；
  \item ValueSSA/MemorySSA tests：锁 SSA conversion、analysis、pass、allocator；
  \item machine tests：锁 machine lowering、select、layout、emit、encode、object/runtime/observe；
  \item compiler driver tests：锁 CLI、模式和输出行为；
  \item extension runtime tests：通过 RISC-V 汇编、链接和 qemu 实际运行扩展语言 witness。
\end{itemize}

每次修复真实 bug 后，都会补对应 regression。例如函数值实参中 \texttt{getint()}
被重复求值的问题，就通过 runtime witness 锁定：同一个普通实参必须只求值一次，
然后根据 runtime tag 选择正确 callable 分支。

\subsection{Verifier 与 dump}

项目大量使用 verifier 和 dump expectation。Verifier 负责检查结构契约，dump
则让 IR/SSA/machine artifact 可观察。它们的价值在于：很多 bug 不是最终程序崩溃
才暴露，而是在中间层就能被定位。

这也形成了一个开发经验：复杂 feature 应避免依赖黑箱式成功。如果中间表示不可观察，
修复往往只能依赖最终现象反推原因，效率低且容易误判。

\subsection{全量回归}

项目中一个重要纪律是：阶段性改动必须经过 focused test 和 broad test 双重验证。
典型顺序是：

\begin{enumerate}[label=\arabic*.]
  \item 先运行最小 witness；
  \item 再补 runtime 或 regression test；
  \item 再跑相关子套件；
  \item 最后跑 \texttt{make test} 全量测试。
\end{enumerate}

这种流程使得 vibe coding 不至于变成仅凭直觉推进实现：实现可以快速迭代，
但是否接受改动必须由测试结果决定。

\section{经验教训}

\subsection{分层比局部技巧更重要}

很多复杂问题不是靠某一个局部条件分支解决的，而是靠正确分层降低复杂度。
canonical IR、lower IR、ValueSSA、MemorySSA、machine backend 的边界越清晰，
后续扩展和调试就越容易。

\subsection{语言扩展必须保守打开}

函数值和闭包的实现尤其说明这一点。表面上只是“把函数当参数传”，但很快会牵出
local alias、runtime tag、returned payload、closure env、callable-object IR、
exactly-once evaluation 等一整串问题。如果过早宣称完整一等函数，后续实现和验证成本会迅速上升。
项目选择了保守 slice：能支持的 family 稳定支持，不能支持的 family 明确 gate。

\subsection{性能优化必须讲证据}

\texttt{-perf} 线也带来一个教训：优化不能只看静态代码形状。某个 transform 是否值得保留，
要看真实热点、runtime A/B、compile-time 成本和正确性风险。有些优化在一个 witness 上
效果很好，在另一个 witness 上却可能带来负收益。

\subsection{测试不是收尾，是设计的一部分}

项目越往后，测试越不是实现后的附属品，而是 feature contract 的表达方式。
很多“当前到底支持什么”的答案，不只在代码里，也在测试矩阵里。

\section{总结}

本项目最终形成了一个三模式编译器框架：\texttt{-riscv} 负责稳定基础编译与
RISC-V 输出，\texttt{-perf} 负责中端优化与 family-aware allocator 实验，
\texttt{-extension} 负责语言特性扩展与 callable-object IR 探索。三个模式共享
分层 IR、SSA、后端和测试基础设施，但各自有清晰边界。

从课程项目角度看，它的主要价值不只在于“实现了多少功能”，更在于形成了一套可以继续
扩展的工程方法：先分层，再验证；先保守支持，再逐步推广；先让中间状态可观察，再谈
优化和复杂语义。本项目的核心经验是：面对复杂编译器行为时，不应急于扩大功能表面，
而应优先建立 verifier、测试矩阵和可观察的中间层契约。

\section*{PS}
\addcontentsline{toc}{section}{PS}

\begin{enumerate}[label=\arabic*.]
  \item 想要做好这个真的好难，最花时间的一集。
  \item 已经尽力在让报告写得像人话了！但是真的很难表达明白 www。
  \item OJ 真的什么报错都不给啊喂！只能对着几十万行的项目发呆 debug 是人类能战胜的嘛……
  （\texttt{lv8 hanoi2.sy} 的 RE 陪伴到我做完，还是不知道为什么会 RE。）
  \item vibe 有风险，管理需谨慎。一晚上 AI 写的代码能让我彻底不认识这个项目。
\end{enumerate}

\appendix

\section{Showcase 源程序与复现命令}

\subsection{主 showcase 源程序}

\begin{lstlisting}[language=C]
struct Box { pair p; int bias; };

int sum_box(struct Box b) {
  return b.p.first + b.p.second + b.bias;
}
struct Box mkbox(int a, int b, int c) {
  struct Box box = {{a, b}, c};
  return box;
}

int apply(int f(int), int x) { return f(x); }
int pass(int h(int f(int), int x), int f(int), int x) {
  return h(f, x);
}
int make(int base)(int) {
  return closure [base] int (int y) { return base + y; };
}
int main() {
  struct Box b = mkbox(1, 2, 3);
  int f(int) = make(sum_box(b));
  return pass(apply, f, 4);
}
\end{lstlisting}

\subsection{defer showcase 源程序}

\begin{lstlisting}[language=C]
int main() {
  int x = 1;
  defer { x = x + 10; }
  capdefer(y = x) putint(y);
  x = 2;
  fndefer putint(x);
  return x;
}
\end{lstlisting}

\subsection{复现命令}

\begin{lstlisting}[language=bash]
build/tools/dump_middle_stage --extension ir showcase.sy
build/tools/dump_middle_stage --extension lower-ir showcase.sy
build/compiler -extension showcase.sy -o showcase.s

build/tools/dump_middle_stage --extension ir defer_showcase.sy
build/compiler -extension defer_showcase.sy -o defer_showcase.s
\end{lstlisting}

正文中的 dump 片段即由上述命令生成。报告正文只保留关键片段，是为了突出设计点并控制篇幅。

\section{项目规模统计}

下表按源码仓库中的主要目录统计行数。统计口径为排除以下目录后的文本行数：

\begin{lstlisting}[language={}]
build/
asanbuild/
third_party/
tmp/
\end{lstlisting}

\texttt{lesson/} 也计入项目产物，因为它承担了设计复盘、代码阅读和后续开发记忆的角色。

\begin{center}
\begin{tabular}{lr}
\toprule
部分 & 行数 \\
\midrule
\texttt{src/} & 235,217 \\
\texttt{include/} & 17,939 \\
\texttt{src + include} & 253,156 \\
\texttt{tests/} & 238,529 \\
\texttt{docs/} & 40,349 \\
\texttt{lesson/} & 83,552 \\
\texttt{reports/} & 4,983 \\
\texttt{tools/} & 7,034 \\
\midrule
项目总量 & 629,432 \\
\bottomrule
\end{tabular}
\end{center}

这组数字不用于证明“代码越多越好”。它更适合说明项目的工程重心：实现与测试规模基本同量级，
而 \texttt{docs/} 与 \texttt{lesson/} 也占了相当比例。对一个快速迭代的编译器项目来说，
这些文字产物不是装饰，而是防止后续 agent 和人类读者迷路的地图。

\section{运行与测试工具链}

本附录给出从“写一个小程序”到“跑完整测试矩阵”的常用命令。以下命令默认在仓库根目录执行。

\subsection{构建编译器}

\begin{lstlisting}[language=bash]
make
make compiler
make tools
\end{lstlisting}

默认 \texttt{make} 只构建 \texttt{build/compiler}，不会自动跑全量测试；这是为了兼容课程工具链中
“先 build 编译器，再由外部评测运行”的习惯。\texttt{make tools} 会构建 dump/profile/allocator
诊断工具。

\subsection{编译一个源程序}

\begin{lstlisting}[language=bash]
build/compiler -riscv input.sy -o output.s
build/compiler -perf input.sy -o output.s
build/compiler -extension input.sy -o output.s

build/compiler --enforce-all-paths-return-check -riscv input.sy -o output.s
\end{lstlisting}

三种模式分别对应基础 RISC-V 输出、性能优化路径和语言扩展路径。strict return gate 可以叠加在
任意模式前，用来开启第 7 章介绍的 all-paths-return 审计。

\subsection{从汇编链接并运行}

生成的 \texttt{.s} 可以手动通过 RISC-V 工具链链接运行：

\begin{lstlisting}[language=bash]
clang output.s -c -o output.o \
  -target riscv32-unknown-linux-elf -march=rv32im -mabi=ilp32

ld.lld output.o -L/opt/lib/riscv32 -lsysy -o output
qemu-riscv32-static ./output
\end{lstlisting}

仓库也提供了 \texttt{make compile-one} 与 \texttt{make run-one} 封装这条链路：

\begin{lstlisting}[language=bash]
make compile-one SRC=input.sy MODE=riscv
make run-one SRC=input.sy MODE=riscv
make run-one SRC=input.sy MODE=extension IN=input.txt RUNOUT=stdout.txt
make run-one SRC=input.sy MODE=perf RETURN_CHECK=1
\end{lstlisting}

\subsection{dump 中间层}

调试编译器时，最常用的是把不同阶段 dump 出来比较：

\begin{lstlisting}[language=bash]
build/tools/dump_middle_stage --extension ir input.sy
build/tools/dump_middle_stage --extension lower-ir input.sy
build/tools/dump_middle_stage --extension value-ssa-default input.sy

build/tools/dump_machine_stage machine-ir input.sy
build/tools/dump_machine_stage select input.sy
build/tools/dump_machine_stage emit input.sy

build/tools/dump_alloc_stage initial input.sy
build/tools/dump_alloc_stage final input.sy
\end{lstlisting}

\texttt{dump\_middle\_stage} 主要看 IR/lower IR/SSA；\texttt{dump\_machine\_stage} 主要看后端
machine-facing artifact；\texttt{dump\_alloc\_stage} 用于观察 allocator 计划和最终分配结果。

\subsection{测试矩阵}

全量测试入口是：

\begin{lstlisting}[language=bash]
make test
make check
\end{lstlisting}

常用分项测试包括：

\begin{lstlisting}[language=bash]
make test-lexer-regression
make test-parser-regression
make test-semantic-regression
make test-ir-regression
make test-lower-ir-regression
make test-value-ssa-regression
make test-value-ssa-alloc
make test-memory-ssa-pass
make test-machine-ir
make test-machine-select
make test-compiler-driver
make test-extension-runtime
\end{lstlisting}

此外还有更严格的构建/检查模式：

\begin{lstlisting}[language=bash]
make test-fanalyzer
make test-asan
make test-strict-warnings
\end{lstlisting}

\subsection{extension runtime 与 SysY suite}

\texttt{make test-extension-runtime} 会调用：

\begin{lstlisting}[language=bash]
tools/test_extension_runtime.sh
\end{lstlisting}

这条脚本会把 \texttt{-extension} witness 编译成 RISC-V 汇编，再用下面的工具真正运行，检查
stdout 和 exit status。

\begin{lstlisting}[language={}]
clang
ld.lld
qemu-riscv32-static
\end{lstlisting}

对于外部 SysY 测试集，可以使用 suite sweep 工具：

\begin{lstlisting}[language=bash]
python3 tools/sweep_sysy_suite.py path/to/suite --compiler build/compiler
\end{lstlisting}

如果测试集带真实输入文件，应使用真实输入运行；空 stdin 的 qemu 结果不能作为性能或正确性结论。

\subsection{性能与诊断工具}

\begin{lstlisting}[language=bash]
python3 tools/profile_compile_pipeline.py input.sy
build/tools/profile_compiler_stages input.sy
build/tools/profile_backend_layers input.sy
build/tools/diag_allocator_stages input.sy

python3 tools/analyze_riscv_asm.py output.s
python3 tools/classify_machine_select_hotspots.py dump.txt
\end{lstlisting}

\texttt{-perf} 相关结论应优先使用真实 runtime A/B，而不是只看静态指令数或“汇编看起来更短”。
这些诊断工具主要用于定位热点、解释形状和缩小问题范围。

\subsection{报告编译}

本报告使用 XeLaTeX 编译：

\begin{lstlisting}[language=bash]
latexmk -xelatex -interaction=nonstopmode \
  -halt-on-error -outdir=reports reports/course_report.tex

latexmk -c -outdir=reports reports/course_report.tex
\end{lstlisting}

\section{进一步阅读路线}

报告正文刻意只展开了最有代表性的设计点。若希望继续核对实现细节，可以按下面路线阅读源码、
测试、设计文档和 lesson。每组先给阅读目的，再给入口文件；这些不是完整目录树，只是适合继续
顺藤摸瓜的起点。

\subsection{项目总览与工作记忆}

先看这些文件，可以了解 agent 如何启动、项目当前处在什么阶段，以及哪些工程事实已经被记录为
工作记忆。

\begin{lstlisting}[language={}]
AGENTS.md
docs/NEXT_STEPS.md
docs/ENGINEERING_MEMORY.md
lesson/README.md
\end{lstlisting}

\subsection{前端、语义与 strict return}

这一组用于补足词法、语法、语义和 all-paths-return 的细节，其中控制流摘要文件对应正文第 7 章。

\begin{lstlisting}[language={}]
lesson/core/lexer_lesson.md
lesson/core/parser_lesson.md
lesson/core/semantic_lesson.md
src/parser/
src/semantic/
src/semantic/semantic_core_flow.inc
tests/semantic/semantic_regression_scope_cf.inc
tests/semantic/semantic_regression_callable_flow.inc
tests/compiler/compiler_driver_test.c
\end{lstlisting}

\subsection{IR、lower IR 与可观察 dump}

这一组对应正文里的 canonical IR、lower IR、verifier 和 dump。想核对 showcase 片段，可以从
\texttt{dump\_middle\_stage} 重新生成。

\begin{lstlisting}[language={}]
lesson/core/ir_lesson.md
lesson/core/lower_ir_lesson.md
docs/ir/LOWER_IR_DESIGN.md
include/ir.h
src/ir/
tests/ir/
include/lower_ir.h
src/lower_ir/
tests/lower_ir/
build/tools/dump_middle_stage
\end{lstlisting}

\subsection{SSA、MemorySSA 与优化}

这一组适合继续看 ValueSSA、MemorySSA、优化 pass 和 \texttt{-perf} 工作流。正文只讲了主线，
详细 pass 取舍和 A/B 记录主要在优化计划里。

\begin{lstlisting}[language={}]
lesson/ssa/value_ssa_lesson.md
lesson/ssa/memory_ssa_lesson.md
docs/ssa/VALUE_SSA_DESIGN.md
docs/ssa/MEMORY_SSA_DESIGN.md
src/value_ssa/
src/value_ssa_pass/
tests/value_ssa/
src/memory_ssa/
tests/memory_ssa/
docs/OPTIMIZATION_PLAN.md
docs/GENERAL_OPT_PASSES_TODO.md
\end{lstlisting}

\subsection{family-aware allocator}

正文的 allocator 例子只展示了关键字段。若要继续看 family、priority、targeted eviction
和 result shaping，可以从这些入口读。

\begin{lstlisting}[language={}]
lesson/ssa/value_ssa_alloc_lesson.md
docs/ssa/VALUE_SSA_ALLOCATOR_PLAN.md
docs/ssa/VALUE_SSA_ALLOCATOR_RESULT_SHAPING_PLAN.md
src/value_ssa_alloc/
tests/value_ssa/
\end{lstlisting}

\subsection{后端与 RISC-V 输出}

这一组对应 \texttt{-riscv} 与 machine backend。machine 层文件很多，建议先从 plan 和
lowering 主线读起，再看 object/runtime/observe。

\begin{lstlisting}[language={}]
lesson/machine/README.md
docs/backend/MACHINE_IR_PLAN.md
docs/backend/MACHINE_SELECT_PLAN.md
docs/backend/MACHINE_EMIT_PLAN.md
docs/backend/MACHINE_ENCODE_PLAN.md
include/machine/
src/machine/lowering/
tests/machine/lowering/
src/machine/object/
src/machine/runtime/
src/machine/observe/
\end{lstlisting}

\subsection{\texttt{-extension} 语言扩展}

这一组对应正文第 5 章和第 6 章。建议先读总览，再按 defer、aggregate、function value、
closure、callable-object IR 分组深入。

\begin{lstlisting}[language={}]
lesson/language/defer_family_lesson.md
docs/language/DEFER_FEATURE_PLAN.md
docs/language/FUNCTION_EXIT_DEFER_PLAN.md
docs/language/DEFER_CAPTURE_PLAN.md

lesson/language/aggregate_lesson.md
lesson/language/aggregate_boundary_lesson.md
docs/language/PAIR_FEATURE_PLAN.md
docs/language/STRUCT_FEATURE_PLAN.md
docs/language/AGGREGATE_PARAM_RETURN_PLAN.md

lesson/language/function_values_lesson.md
lesson/language/function_implementation_map_lesson.md
lesson/language/function_supported_code_lesson.md

lesson/language/closure_object_lesson.md
lesson/language/returned_callable_lesson.md
lesson/language/returned_closure_transport_lesson.md
lesson/language/higher_order_callable_lesson.md

lesson/language/callable_object_ir_lesson.md
docs/ir/CALLABLE_OBJECT_IR_DESIGN.md
docs/language/CALLABLE_OBJECT_IR_PLAN.md
\end{lstlisting}

扩展语言最密集的功能和 gate 测试入口在：

\begin{lstlisting}[language={}]
tests/compiler/compiler_driver_test.c
tests/semantic/semantic_regression_test.c
\end{lstlisting}

\subsection{测试体系}

最后，如果要从测试角度理解“项目到底锁住了哪些行为”，可以从这些入口开始。

\begin{lstlisting}[language={}]
lesson/core/tests_lesson.md
tests/lexer/
tests/parser/
tests/semantic/
tests/ir/
tests/lower_ir/
tests/value_ssa/
tests/memory_ssa/
tests/machine/
tests/compiler/
\end{lstlisting}

\end{document}
